\documentclass[11pt,letterpaper]{article}
\usepackage{fontspec}
\usepackage{tocloft}
\usepackage{multicol}
\usepackage{nicefrac}
\usepackage[left=1in,right=1in,top=1in,bottom=1in]{geometry}
\setmainfont[Scale=1.2,AutoFakeBold=1.5,AutoFakeSlant=0.3]{Doves Type}

\widowpenalty=10000
\clubpenalty=10000
\interlinepenalty=500

\title{Chicken Fricassee}
\author{}
\date{}

\begin{document}

\maketitle
\thispagestyle{empty}

\begin{quote}\textit{Bone-in chicken thighs are browned in butter, then braised in a pale roux-thickened sauce with cremini mushrooms, onion, and garlic. Dry white wine and chicken stock build the base; heavy cream finishes the sauce until silky and pale. Serve over egg noodles, mashed potatoes, or rice. Serves 4--6.}\end{quote}

\section*{Ingredients}
\setlength{\columnsep}{20pt}
\begin{multicols}{2}
\noindent
    Chicken thighs, bone-in, skin-on \dotfill 8 (~3~lb.) \\
    Kosher salt and black pepper \dotfill as needed \\
    Unsalted butter \dotfill 4~Tbsp. \\
    Cremini mushrooms \dotfill 10~oz. \\
    Yellow onion \dotfill 2~medium \\
    Garlic cloves \dotfill 3 \\
    Bay leaf \dotfill 1 \\
    Dried thyme \dotfill 1~tsp. \\
    All-purpose flour \dotfill 3~Tbsp. \\
    Dry white wine \dotfill \nicefrac{1}{2}~cup \\
    Chicken stock \dotfill 3~cups \\
    Heavy cream \dotfill \nicefrac{2}{3}~cup \\
    Dried parsley \dotfill 1~tsp. \\
    Nutmeg \dotfill pinch \\
\end{multicols}

\section*{Directions}

\noindent
Pat \textbf{chicken~thighs} dry and season generously with \textbf{salt} and \textbf{black~pepper} ---
Slice \textbf{cremini~mushrooms}; set aside in \textit{Medium Bowl~\#1} (mushrooms) ---
Dice \textbf{yellow~onions} and mince \textbf{garlic}; combine in \textit{Small Bowl~\#1} (aromatics) ---
Measure \textbf{dried~thyme}, \textbf{bay~leaf}, \textbf{flour}, \textbf{white~wine}, \textbf{chicken~stock}, \textbf{heavy~cream}, \textbf{dried~parsley}, and \textbf{nutmeg}

\begin{enumerate}
    \item Heat a large Dutch oven or deep skillet over medium-high heat. Melt 2~Tbsp. \textbf{butter} until foaming subsides. Working in batches if needed to avoid crowding, add \textbf{chicken~thighs} skin-side down and brown for \textit{4--5~minutes} until skin is deep golden and releases easily from the pan. Flip and brown the second side for \textit{3--4~minutes} until lightly golden. Transfer \textbf{chicken} to \textit{Large Bowl~\#1} and set aside.

    \item Reduce heat to medium-high. Add 1~Tbsp. \textbf{butter} to the pot. Add sliced \textbf{cremini~mushrooms} (\textit{Medium Bowl~\#1}) and cook, stirring occasionally, for \textit{12--15~minutes}. They will release a pool of liquid first, then dry out and begin to brown as that moisture evaporates. Continue until mushrooms are deeply browned, the pan looks dry, and edges look crisp at the rim. Transfer cooked \textbf{mushrooms} back to \textit{Medium Bowl~\#1} and set aside.

    \item Reduce heat to medium. Add remaining 1~Tbsp. \textbf{butter}. Add diced \textbf{onion} (\textit{Small Bowl~\#1}), \textbf{bay~leaf}, and 1~tsp. \textbf{dried~thyme}. Cook, stirring often, for \textit{10--12~minutes} until \textbf{onion} is translucent and fully softened, edges are lightly golden, and the pieces feel tender when pressed with a spatula. If \textbf{onion} browns too quickly, lower heat slightly and add a splash of \textbf{chicken~stock} to slow cooking. Add minced \textbf{garlic} (\textit{Small Bowl~\#1}) and cook for \textit{1--2~minutes} more until fragrant without browning.

    \item Sprinkle 3~Tbsp. \textbf{flour} evenly over the \textbf{onion} mixture and stir constantly for \textit{2~minutes} until the roux smells nutty and toasty and no dry flour remains. The paste should look smooth and pale gold. If roux clumps, press and stir vigorously with a wooden spoon until smooth. Reduce heat to medium-low if the roux darkens too fast.

    \item Pour in \nicefrac{1}{2}~cup \textbf{dry~white~wine}, scraping the bottom of the pot to dissolve browned bits. Simmer for \textit{2--3~minutes}, stirring, until the wine reduces slightly and the harsh alcohol smell fades. The liquid should look glossy and lightly thickened at the edges.

    \item Whisk in 3~cups \textbf{chicken~stock} until smooth. Return \textbf{chicken} (\textit{Large Bowl~\#1}) to the pot skin-side up. Bring to a gentle simmer, taste the liquid, and season with \textbf{salt} and \textbf{black~pepper} as needed. Cover and simmer over low heat for \textit{15~minutes}, then uncover and continue simmering for \textit{18--22~minutes} until \textbf{chicken} is fall-apart tender, meat pulls easily from the bone, and the thickest part registers \textit{165°F} on an instant-read thermometer. Continue simmering in \textit{3--5~minute} increments if the center is still pink or firm.

    \item Transfer \textbf{chicken} back to \textit{Large Bowl~\#1} and tent loosely with foil to keep warm. Stir \nicefrac{2}{3}~cup \textbf{heavy~cream}, 1~tsp. \textbf{dried~parsley}, and a pinch of \textbf{nutmeg} into the sauce. Simmer gently over medium-low heat for \textit{3--5~minutes}, stirring, until the sauce is silky, pale ivory, and coats the back of a spoon (a line drawn through it with your finger should hold clearly). Taste and adjust \textbf{salt}, \textbf{black~pepper}, and \textbf{nutmeg} as needed. If the sauce breaks or looks greasy, remove from heat and whisk vigorously until smooth. If too thin, simmer uncovered for another \textit{2--3~minutes}.

    \item Return cooked \textbf{mushrooms} (\textit{Medium Bowl~\#1}) and \textbf{chicken} (\textit{Large Bowl~\#1}) to the sauce. Warm over low heat for \textit{2--3~minutes}, spooning sauce over the \textbf{chicken}, until everything is hot throughout. Discard the \textbf{bay~leaf}. Serve immediately while the sauce is glossy and the \textbf{chicken} is tender.
\end{enumerate}

\newpage

{\footnotesize
\setlength{\columnsep}{20pt}
\setlength{\multicolsep}{6pt}
\begin{multicols}{2}
\setlength{\parindent}{0pt}
\setlength{\parskip}{2pt}
\setlength{\itemsep}{0pt}
\setlength{\parsep}{0pt}

\section*{Equipment Required}
\begin{itemize}
    \item Large Dutch oven or deep skillet with lid (6-quart preferred)
    \item Wooden spoon or spatula
    \item Whisk
    \item Instant-read thermometer
    \item Large prep bowl (1)
    \item Medium prep bowl (1)
    \item Small prep bowl (1)
    \item Cutting board and chef's knife
    \item Measuring cups and spoons
\end{itemize}

\section*{Hints and Notes}

\subsection*{Yield}
\begin{itemize}
    \item Serves 4--6 (8 bone-in thighs with generous sauce).
\end{itemize}

\subsection*{Mise en Place}
\begin{itemize}
    \item \textbf{Small Bowl~\#1} --- aromatics: diced \textbf{yellow~onion} and minced \textbf{garlic} (onion added in step 3 first; garlic added after \textit{10--12~minutes})
    \item \textbf{Medium Bowl~\#1} --- sliced \textbf{cremini~mushrooms}; holds cooked mushrooms after step 2
    \item \textbf{Large Bowl~\#1} --- browned then braised \textbf{chicken~thighs}
    \item Measure 1~tsp. \textbf{dried~thyme}, 1~tsp. \textbf{dried~parsley}, 3~Tbsp. \textbf{flour}, \nicefrac{1}{2}~cup \textbf{white~wine}, 3~cups \textbf{chicken~stock}, \nicefrac{2}{3}~cup \textbf{heavy~cream}, and a pinch of \textbf{nutmeg}
    \item Pat \textbf{chicken} dry before browning; surface moisture inhibits browning and dilutes the fond
\end{itemize}

\subsection*{Ingredient Tips}
\begin{itemize}
    \item \textbf{Chicken~thighs}: bone-in, skin-on thighs braise best and stay moist. Skin-on pieces also render fat that enriches the sauce.
    \item \textbf{Cremini~mushrooms}: deeper flavor than white button mushrooms. Quarter or thickly slice so they brown instead of steaming.
    \item \textbf{Dry~white~wine}: Pinot Grigio, Sauvignon Blanc, or unoaked Chardonnay add acidity without sweetness. Avoid sweet or oaked wines that muddy the pale sauce.
    \item \textbf{Dried~herbs}: 1~tsp. \textbf{dried~thyme} replaces about 3 fresh sprigs. \textbf{Dried~parsley} is milder than fresh; add at the finish so its flavor stays in the sauce.
    \item \textbf{Chicken~stock}: low-sodium stock gives more control over salting. Homemade stock deepens flavor further.
\end{itemize}

\subsection*{Preparation Tips}
\begin{itemize}
    \item \textbf{Browning}: do not crowd the pot. Crowded \textbf{chicken} steams instead of browning and leaves a thin fond. Work in two batches if your Dutch oven is smaller than 6-quart.
    \item \textbf{Mushrooms}: expect a wet phase before browning begins. Cook the full \textit{12--15~minutes} until the pan is dry and the mushrooms are deeply browned. Wet mushrooms make a watery sauce when returned at the end.
    \item \textbf{Onions}: give diced \textbf{onion} the full \textit{10--12~minutes} to soften and sweeten before adding \textbf{garlic}. Adding \textbf{garlic} earlier risks burning it during the longer sweat.
    \item \textbf{Roux}: cook the flour fully before adding liquid. A raw-flour taste means the roux needed more time; continue stirring over medium-low heat until nutty.
    \item \textbf{Braise}: keep the simmer gentle after uncovering. A rolling boil tightens \textbf{chicken} and can break the cream-finished sauce.
    \item \textbf{Cream~finish}: add \textbf{heavy~cream} after the \textbf{chicken} is cooked through. High heat after dairy goes in can cause curdling.
\end{itemize}

\subsection*{Make Ahead \& Storage}
\begin{itemize}
    \item Flavor improves after an overnight rest in the refrigerator. Cool completely, then refrigerate up to \textit{3~days}.
    \item Reheat gently on the stovetop over low heat, stirring and adding a splash of \textbf{chicken~stock} or \textbf{cream} if the sauce has thickened.
    \item Freeze up to \textit{2~months} without the \textbf{cream} finish for best texture; stir in \textbf{cream} when reheating.
    \item Peak moment is right off the stove, when the sauce is glossy and \textbf{chicken} is fall-apart tender.
\end{itemize}

\subsection*{Serving Suggestions}
\begin{itemize}
    \item Serve over egg noodles, mashed potatoes, or rice.
    \item Spoon the sauce generously over the \textbf{chicken} and mushrooms at the table.
    \item Crusty bread works well for soaking up any sauce left in the bowl.
\end{itemize}

\end{multicols}
}

\end{document}
